\documentclass[11pt]{article}
% Supplementary material for:
% "The Anchor Accumulation Problem: Bounded Anchor Sets, a Soft--Hard
%  Protection Law, and the Auditable Cost of Regularized Continual
%  Learning"
% Appendix A: statistics rationale, hyperparameters, reproducibility.
% Appendix B: full per-task accuracy matrices (Permuted-MNIST).

\usepackage[margin=1in]{geometry}
\usepackage{amsmath,amssymb}
\usepackage{booktabs}
\usepackage{graphicx}
\usepackage{url}
\graphicspath{{figs/}}

\newcommand{\delp}{DeLP}
\newcommand{\acc}{\mathrm{acc}}

\title{Supplementary Material:\\
The Anchor Accumulation Problem: Bounded Anchor Sets, a Soft--Hard
Protection Law, and the Auditable Cost of Regularized Continual
Learning}
\author{Ranxi Zhou}
\date{}

\begin{document}
\maketitle

\appendix

% ============================================================
\section{Hyperparameters and reproducibility}
% ============================================================
\label{appx:hyper}

\paragraph{Statistics.}
All comparisons between two arms are \emph{paired by seed}: each seed
supplies one run per arm under an identical protocol, and every reported
$\Delta$ is the per-seed difference. We report the paired $t$ statistic
$t(9)$ together with the two-sided p-value of the \emph{Wilcoxon
signed-rank} test on the same paired differences, and we report the
Shapiro--Wilk p-value on those differences so that the normality
assumption underlying the $t$ test is visible rather than implicit.
The choice to report both is deliberate. Whether a benchmark's headline
number supports a claim depends on the variability of the metric under
resampling, not only on its mean. Several authors have argued this point
for generative evaluation~\cite{lucic2018,melis2018}, for metric
learning~\cite{musgrave2020}, for deep reinforcement
learning~\cite{agarwal2021}, and, closest to our setting, for
reproducibility in machine-learning scholarship as a
whole~\cite{lipton2019}. A single significance test applied to ten
seeds is exactly the setting in which those critiques bite: with ten
paired samples the parametric $t$ test has assumptions that the data
may not satisfy, while the distribution-free signed-rank test is
valid at this sample size by construction. We therefore report both
and let the reader see where they diverge.
The two tests agree on every headline comparison: the FIFO budget beats
unbounded at $M{=}2$ and $M{=}5$, $T{=}10$ and $T{=}20$, with
Wilcoxon $p{=}0.002$ (the smallest attainable two-sided value at
$n{=}10$ paired samples) in all four cases. One case deserves the
explicit note: at $M{=}5$, $T{=}10$ the paired differences depart
from normality (Shapiro--Wilk $p{=}0.024$), so the $t$ test there
rests on an assumption the data do not satisfy, whereas the
distribution-free Wilcoxon result ($p{=}0.002$) does not. We therefore
treat the Wilcoxon value as the primary significance statement
throughout, and the $t$ statistic as a secondary, parametric
cross-check. The significance threshold is $\alpha{=}0.05$ throughout,
with no multiplicity correction: the policy$\times$budget grid of the
Split-CIFAR-100 table in the main text contains twelve cells and is
read as a descriptive map of the policy landscape, with per-cell
$p$-values reported so that the reader can judge each claim
individually. All $p$-values are recomputed from the raw per-seed
records by \texttt{audit/audit\_nonparametric.py}, which is included in
the accompanying package.

\paragraph{Permuted-MNIST / Fashion-MNIST (main-text experiment
sections).}
Architecture as in the audit~\cite{zhou2026delp}: a 784--256 fully
connected trunk whose pre-activation is modulated by $K{=}T$ dendritic
segments (linear context-to-segment map, winning-segment readout),
hidden activity sparsified by $k$-WTA with $k{=}64 = N/4$, linear
10-way head. Adam, lr $10^{-3}$, batch 256, two epochs and 10{,}000
samples per task; permutation matrices drawn once per seed and shared
by all arms. Diagonal Fisher and usage statistics from 300 and 500
samples per task respectively (max-log-probability Fisher variant);
$\lambda_{\delp}{=}5000$, corrected-SI reference $\lambda{=}100$.
Paired arms use seeds 0--9; the policy grid seeds 0--2 (policy grid
table in the main text); the $\lambda$ grid as in the $\lambda$-interaction
table of the main text. Every run seeds
\texttt{torch.manual\_seed} before model construction. Code and
per-seed JSON records (including \texttt{survivor\_history} and
\texttt{acc\_matrix}) are released with the repository.

\begin{sloppypar}
\paragraph{Split-CIFAR-100 (main text, Sec.~on generalization).}
ResNet-18 adapted for $32{\times}32$ inputs ($3{\times}3$ stem,
stride 1, no max-pool) with all BatchNorm layers replaced by GroupNorm
(32 groups). SGD, momentum 0.9, weight decay $5{\times}10^{-4}$, lr
$0.05$, batch 128, five epochs per task. Masked cross-entropy over the
task's $100/T$ classes; task-incremental evaluation over the same classes.
EWC penalty over backbone parameters only, $\lambda{=}0.1$; diagonal
Fisher estimated from 50 single-sample draws of the squared gradient
of the max log-softmax on the finished task, rescaled to unit global
mean; gradient clipping at max norm 5. Tasks are contiguous
$100/T$-class blocks of a fixed random permutation of the 100 classes
(generator seed $100{+}\mathrm{seed}$, matching the released code);
seeds $0$--$9$ control init,
shuffle, and eviction randomness. Per-job crash-safe checkpoints and
the full JSON records (including \texttt{survivor\_history} and
\texttt{acc\_matrix}) are released with the code.
\end{sloppypar}

% ============================================================
\section{Full per-task accuracy matrices}
% ============================================================
\label{appx:matrices}

\begin{sloppypar}
Final-phase accuracy (\%) per task, Permuted-MNIST, $\lambda{=}5000$,
$M{=}5$: mean over seeds ($n{=}10$ for unbounded and fifo, $n{=}3$ for
random/redun/merge). These underlie Fig.~3a and the grandfathering
analysis of the main text.
\end{sloppypar}

\begin{table}[h]
\centering\small
\caption{$T{=}10$.}
\label{tab:appx-t10}
\input{appx_mat_t10.tex}
\end{table}

\begin{table}[h]
\centering\small
\caption{$T{=}20$.}
\label{tab:appx-t20}
\input{appx_mat_t20.tex}
\end{table}

\section*{Data availability}
All code, experiment scripts, and the complete JSON records of every
run (including per-phase accuracy matrices and anchor-stack survivor
histories) are publicly available at
\url{https://github.com/zhouranxi/delp-anchor-budget}. This manuscript
and its full reproducibility package --- including the audit scripts
that re-derive every quoted number from those records --- are archived
on Zenodo (doi:10.5281/zenodo.23668522). The precursor audit study is
archived at doi:10.5281/zenodo.23572226 and the companion
segment-specialization audit at doi:10.5281/zenodo.23637972; the
underlying dataset and code release is
doi:10.5281/zenodo.23639561.

\bibliographystyle{plain}
\bibliography{refs}

\end{document}
